
abéliens, de la
façon suivante : si \(A\in\operatorname{Ob}(\mathcal C)\), on associe à \(A\) l'ensemble des séries formelles \(\psi\in A[[t]]\), vérifiant les conditions suivantes : si \(\psi=a_0+a_1t+\ldots+a_nt^n+\ldots\), alors
\begin{equation}\tag{1.1.1}
  a_0=1;\qquad \forall i\geq 1,\quad a_i\in\Fil^i A .
\end{equation}
Cet ensemble sera noté \(\Ghat_K(A)\), ou \(\Ghat(A)\) si aucune confusion n'est possible ; il est clair que c'est un groupe abélien pour la multiplication des séries formelles, et qu'on définit bien ainsi un foncteur
\[
  \Ghat_K:\underline{K\text{-Alg.fil.}}\longrightarrow \underline{\mathrm{Ab}}.
\]
Si \(A\) est muni de la filtration grossière, \(\Ghat(A)\) est le groupe des séries formelles d'augmentation \(1\), qu'on notera aussi \(1+A[[t]]^+\).

\textbf{Lemme 1.2.} \uline{Le foncteur} \(\Ghat_K\) \uline{est représentable}.

Soit \((X_i)_{i\in N,\ i\geq 1}\), une famille d'indéterminées ; on attribue le poids \(i\) à \(X_i\), et on filtre l'anneau \(K[X_i]_{i\geq 1}\) par le poids par rapport aux \(X_i\). On définit alors une bijection~:
\[
  \Ghat_K(A)\xrightarrow{\sim}
  \Hom_{K\text{-Alg.fil.}}\bigl(K[X_i]_{i\geq 1},A\bigr)
\]
en associant à toute série \(\psi=1+a_1t+\ldots+a_nt^n+\ldots\in\Ghat_K(A)\) l'homomorphisme \(f_\psi\) déterminé par \(f_\psi(X_i)=a_i\). Cet homomorphisme est compatible aux filtrations, et on obtient ainsi un isomorphisme de foncteurs.

On pose \(\phi=1+X_1t+\ldots+X_nt^n+\ldots\)~; pour toute \(K\)-algèbre filtrée \(A\) et toute série \(\psi\in\Ghat(A)\), il existe un homomorphisme et un
seul 
